\documentclass[11pt]{article}

\usepackage[margin=1in]{geometry}
\usepackage{amsmath,amssymb,amsthm,mathtools}
\usepackage{enumitem}
\usepackage{booktabs}
\usepackage{tabularx}
\usepackage{microtype}
\usepackage[hidelinks]{hyperref}
\usepackage{xurl}

\newtheorem{definition}{Definition}
\newtheorem{lemma}{Lemma}
\newtheorem{remark}{Remark}

\title{Challenge Answer Escalation Ladders for Successor Maintenance:\\[0.5ex]
\large Ordered Brief-to-Audit Reveal Policies from Publication Profiles and Disclosure Packets}
\author{}
\date{Unified archive note v0.06 (2026-03-21; archive v479)}

\begin{document}
\maketitle

\begin{abstract}
Once answer publication profiles and disclosure packets exist, later brief papers still face one repeated policy choice under space pressure: in what maintained order should one reader move from the terse printed answer to richer support? Some notes jump straight from a brief sentence to a full audit trail; others point to a disclosure packet without ever naming the exact shipped carrier slot. This note gives that reveal order one home. It defines a compact \emph{challenge answer escalation ladder}: one audience/question policy object ordering the already-maintained brief-inline, expanded-inline, exact-location, full-audit, and request-packet companions for the same emitted answer.
\end{abstract}

\paragraph{Non-redundancy note.}
Synthesis~50 owns smallest-sufficient answer citation dockets, Synthesis~51 owns inline publication-budget choices, Synthesis~52 owns exact public-on-request disclosure packets, Synthesis~49 owns full leftward answer-audit trails, Synthesis~48 owns exact answer carrier slots, Synthesis~31 owns the downstream suffix, Synthesis~32 owns the stable release spine, and Synthesis~17 owns the concrete worked instantiation.\cite{series:challengeanswercitationdockets,series:challengeanswerpublicationprofiles,series:challengeanswerdisclosurepackets,series:challengeansweraudittrails,series:challengeanswercarrierslots,series:downstreamnormalforms,series:releasespine,series:workedexample}
This note owns only the ordered reveal policy among those already-maintained companions.

\paragraph{Scope.}
This is not a new answer card, carrier slot, audit trail, citation docket, publication profile, or disclosure packet.
It is a narrow policy note about ordering those imported objects so one terse paper can reveal more support in a maintained way rather than improvising that escalation sequence at every reuse site.

\paragraph{Imports.}
Answer citation dockets are imported from Synthesis~50.\cite{series:challengeanswercitationdockets}
Answer publication profiles are imported from Synthesis~51.\cite{series:challengeanswerpublicationprofiles}
Answer disclosure packets are imported from Synthesis~52.\cite{series:challengeanswerdisclosurepackets}
The concrete worked instantiation is imported from Synthesis~17.\cite{series:workedexample}

\paragraph{Exports.}
This note exports (i) challenge answer escalation ladders, (ii) a monotone-reveal rule, (iii) an imported-consistency rule, (iv) a same-answer identity rule, and (v) a shared-escalation-vocabulary rule saying that later papers may cite one ladder only when the live question is the exact brief-to-audit reveal order for one audience/question answer and may mirror its repeated step labels only as imported ladder vocabulary. Otherwise later terse papers should default to the answer-side terminal citation normal form in Synthesis~54 together with the paired terminal discipline and paired cutover rule in Synthesis~74--75. Choosing the smallest sufficient handoff object for one concrete referee follow-up class belongs one note later in Synthesis~54.\cite{series:challengeanswerreviewmenus,series:pairedterminalcitationdiscipline,series:pairedterminalcutoverrules}

\section{Why answer escalation ladders deserve their own note}
Publication profiles already tell a terse paper what to print by default and what a larger inline form looks like.
Citation dockets already tell it which object to cite if the reader asks for exact location or full audit.
Disclosure packets already tell it what exact archive slice should be handed over on request.
But the series still lacks one maintained answer to the practical referee question: \emph{in what order should these already-maintained forms be revealed?}
Without one ordered policy object, later notes either skip intermediate forms or silently reinvent them.
So the archive benefits from one thin ladder whose only job is to keep the reveal sequence stable.

\begin{definition}[Challenge answer escalation ladder]
Fix one concrete base$\to$successor comparison, one audience key $a$, and one downstream question key $q$.
A \emph{challenge answer escalation ladder} is a tuple
\[
L(a,q)=(a,q,P,C,D,E_1,E_2,E_3,E_4,E_5)
\]
containing one imported answer-publication-profile key $P$, one imported answer-citation-docket key $C$, one imported answer-disclosure-packet key $D$, and an ordered five-step list $(E_1,\dots,E_5)$ consisting of a default-inline entry, an expanded-inline entry, an exact-location entry, a full-audit entry, and a request-packet entry for the same audience/question answer.
\end{definition}

\begin{definition}[Monotone-reveal rule]
A challenge answer escalation ladder satisfies the \emph{monotone-reveal rule} if its ordered steps reveal weakly more support in the fixed order
\[
\textsf{default-inline} \prec \textsf{expanded-inline} \prec \textsf{exact-location} \prec \textsf{full-audit} \prec \textsf{request-packet}.
\]
So later notes may move rightward when the reader asks for more support, but they may not reorder those reveal modes ad hoc.
\end{definition}

\begin{definition}[Imported-consistency rule]
A challenge answer escalation ladder satisfies the \emph{imported-consistency rule} if its first two steps agree exactly with the imported publication profile, its exact-location and full-audit steps agree exactly with the imported citation docket, and its final request-packet step agrees exactly with the imported disclosure packet.
So the ladder owns only order, not the content of any rung.
\end{definition}

\begin{definition}[Same-answer identity rule]
A challenge answer escalation ladder satisfies the \emph{same-answer identity rule} if every rung refers to the same audience/question pair and every answer-side rung refers to the same emitted answer card whenever that rung still lives on the answer side of the series.
So the ladder may not switch answers while pretending merely to reveal more support.
\end{definition}

\begin{definition}[Shared-escalation-vocabulary rule]
A challenge answer escalation ladder satisfies the \emph{shared-escalation-vocabulary rule} if its repeated step labels \textsf{default-inline}, \textsf{expanded-inline}, \textsf{exact-location}, \textsf{full-audit}, and \textsf{request-packet} may recur in later review menus or series-spine mirrors only as imported ladder vocabulary: those later objects may name the same ordered reveal ladder, but they do not thereby become the owner of the concrete reveal order for that audience/question answer.
So later terse notes may reuse the five rung names without silently re-owning the reveal policy that the ladder fixed.
\end{definition}

\begin{lemma}[Escalation ladders stabilize reveal policy without changing answer ownership]
Assume every challenge answer escalation ladder satisfies the monotone-reveal rule, the imported-consistency rule, the same-answer identity rule, and the shared-escalation-vocabulary rule.
Then a later paper may cite one escalation ladder as the maintained reveal policy for one audience/question answer without changing the ownership of the emitted answer, exact carrier location, full audit walk, or requestable archive slice.
\end{lemma}

\begin{proof}
The monotone-reveal rule fixes the reveal order, so later notes do not improvise a new sequence under space pressure.
The imported-consistency rule prevents the ladder from rewriting any rung: it must copy the already-maintained publication-profile, citation-docket, and disclosure-packet entries exactly.
The same-answer identity rule prevents the ordered list from drifting across audience/question pairs or across emitted answers.
The shared-escalation-vocabulary rule adds the later-seam guarantee: even when later review menus or series-spine mirrors repeat the same five rung labels, those mirrors still import reveal-order vocabulary from the ladder rather than re-owning the reveal order for that answer.
So the ladder changes only policy packaging: it stabilizes the order in which existing support forms are revealed, but it does not alter what any rung means or who owns it.\qedhere
\end{proof}

\begin{table}[t]
\centering
\small
\begin{tabularx}{\linewidth}{@{}l l X@{}}
\toprule
Rung & Canonical object & Reader question answered \\
\midrule
1 & publication profile (default) & What is the smallest inline form the brief paper should print? \\
2 & publication profile (expanded) & What larger inline provenance bundle may the same paper print without changing the answer? \\
3 & citation docket / carrier slot & Where exactly does that emitted answer live once the reader asks for location? \\
4 & citation docket / audit trail & What full leftward walk discharges the answer once exact location is not enough? \\
5 & disclosure packet & What exact archive slice and validator companions should be handed over on request? \\
\bottomrule
\end{tabularx}
\caption{An escalation ladder orders already-owned late-stage answer companions; it does not create a new answer stage.}
\label{tab:answerescalationladder}
\end{table}

\section{Worked example pointer}
The maintained worked bundle now ships \nolinkurl{example_successor_challenge_answer_escalation_ladders.json}.\cite{series:workedexample}
It records four escalation ladders for the canned Case-B successor.
Concrete keys already in the ledger include \nolinkurl{release_note_public_status_sentence_answer_escalation_ladder} and \nolinkurl{auditor_compact_diff_summary_answer_escalation_ladder}.
For the auditor compact-diff answer, the ladder is now concrete enough to read as one maintained widening discipline rather than as a generic five-rung taxonomy: \textsf{default-inline} stops at the default-inline answer card with only \nolinkurl{base\_release\_id}, \nolinkurl{successor\_release\_id}, and \nolinkurl{compare\_report\_id}; \textsf{expanded-inline} stays on that same emitted answer card but widens the visible provenance bundle to include \nolinkurl{object\_digests}; \textsf{exact-location} stops at carrier slot \nolinkurl{auditor_compact_diff_summary_answer_carrier_slot} in \nolinkurl{example_successor_clause_pack.json}; \textsf{full-audit} stops at audit trail \nolinkurl{auditor_compact_diff_summary_answer_audit_trail}; and \textsf{request-packet} stops at disclosure packet \nolinkurl{auditor_compact_diff_summary_answer_disclosure_packet}. The widening triggers are equally explicit: stay on the default rung when the compact sentence itself is enough, widen to the expanded rung only when more visible inline provenance is needed, widen to the carrier slot only when the referee asks where the shipped clause lives, widen to the audit trail only when the clause location is no longer enough because the semantic pieces or terminal owners are contested, and widen to the disclosure packet only when the bounded archive slice itself must be handed over.

\begin{table}[t]
\centering
\small
\begin{tabularx}{\linewidth}{@{}l l X@{}}
\toprule
Rung & Worked compact-diff object & Widen only if \\
\midrule
\textsf{default-inline} & answer card under profile \nolinkurl{floor} & the terse compact diff no longer carries enough visible provenance \\
\textsf{expanded-inline} & same answer card under profile \nolinkurl{capsule} & exact shipped clause location is requested \\
\textsf{exact-location} & \nolinkurl{auditor_compact_diff_summary_answer_carrier_slot} & the semantic pieces or terminal owners must be audited \\
\textsf{full-audit} & \nolinkurl{auditor_compact_diff_summary_answer_audit_trail} & the bounded public-on-request archive slice itself must be shipped \\
\textsf{request-packet} & \nolinkurl{auditor_compact_diff_summary_answer_disclosure_packet} & do not widen inside the ladder; later follow-up-class choice belongs to the review menu \\
\bottomrule
\end{tabularx}
\caption{One concrete worked escalation ladder for the shipped auditor compact-diff answer. The ladder owns only reveal order, not the later referee-class choice of which imported object to hand over.}
\label{tab:workedcompactdiffescalationladder}
\end{table}

Each ordered step now also exports an explicit \nolinkurl{escalation_step_vocabulary_mode} posture tag, and the worked route/spine surfaces mirror that through a shared \nolinkurl{answer_escalation_step_mode_map} so the repeated five rung labels stay visibly ladder-owned rather than being silently flattened into later review or bridge prose. The worked question-routing catalog now mirrors that escalation-ladder key beside the route's checked sentence, publication-budget map, and review-menu key, so later terse notes can recover the maintained reveal order without detouring through the review-menu import list first. But unless reveal order itself is the live question, later terse papers should now stop at the answer-side terminal citation normal form and paired cutover rule instead of restating the five-rung ladder locally.\cite{series:challengeanswerreviewmenus,series:pairedterminalcitationdiscipline,series:pairedterminalcutoverrules} That makes the worked example more referee-useful under space pressure: the archive now says not only what each rung is, but also exactly where the auditor compact-diff route should stop and what concrete question forces the next widening.

\paragraph{A minimal public answer-escalation-ladder card.}
\begin{table}[t]
\centering
\small
\begin{tabularx}{\linewidth}{@{}l l X@{}}
\toprule
Worked ladder & Ordered public reveal spine & Honest stop rule \\
\midrule
\nolinkurl{auditor_compact_diff_summary_answer_escalation_ladder} & \textsf{default-inline} $\to$ \textsf{expanded-inline} $\to$ \textsf{exact-location} $\to$ \textsf{full-audit} $\to$ \textsf{request-packet} & Stay on the current rung until the live question asks for the next larger maintained support form; do not skip straight from the terse sentence to a full audit or request packet while carrier-slot locality still answers the question \\
\nolinkurl{release_note_public_status_sentence_answer_escalation_ladder} & Same five-rung reveal order on the public-status sentence route & Keep the same emitted answer card while only inline provenance expands; once the live question becomes clause-location, leftward audit, or exact omitted-support packet identity, widen exactly one rung at a time rather than re-owning the ladder locally \\
\bottomrule
\end{tabularx}
\caption{A minimal public answer-escalation-ladder card. This is the smallest answer-side public answer later terse papers can quote directly when the live issue is the maintained reveal order itself rather than the narrower imported rung owner or the later referee-class menu choice.}
\label{tab:minimal-public-answer-escalation-ladder-card}
\end{table}

This card is intentionally non-decorative: it pins the actual worked ladder keys, the five imported rung labels they own, and the honest stop rule that keeps later terse papers from silently jumping over exact-location or full-audit when those are still the first sufficient answer-side owners.\cite{series:workedexample}

\paragraph{A forced public answer-escalation-ladder matrix.}
\begin{table}[t]
\centering
\small
\begin{tabularx}{\linewidth}{@{}l X@{}}
\toprule
Force condition & What is still mechanically forced \\
\midrule
Auditor compact-diff ladder floor & The pair \nolinkurl{auditor_compact_diff_summary_answer_escalation_ladder} plus \nolinkurl{default_inline} still forces the checked answer card \nolinkurl{auditor_compact_diff_summary_answer_card} under inline profile \nolinkurl{floor}; the same ladder plus \nolinkurl{expanded_inline} still forces that same emitted answer card under inline profile \nolinkurl{capsule}; the same ladder plus \nolinkurl{exact_location}, \nolinkurl{full_audit}, or \nolinkurl{request_packet} still forces \nolinkurl{auditor_compact_diff_summary_answer_carrier_slot}, \nolinkurl{auditor_compact_diff_summary_answer_audit_trail}, or \nolinkurl{auditor_compact_diff_summary_answer_disclosure_packet} respectively \\
Release-note status ladder floor & The pair \nolinkurl{release_note_public_status_sentence_answer_escalation_ladder} plus any one rung key still forces the same maintained rung object only while the audience/question pair, emitted-answer lineage, imported publication-profile key, imported citation-docket key, and imported disclosure-packet key all continue to agree with that ladder \\
Staleness trigger floor & This matrix goes stale as soon as the audience/question pair changes, the emitted answer card or sentence-lock lineage changes, a different rung key becomes live, the selected rung object no longer matches the ladder entry, or the ladder ceases to import the same publication-profile / citation-docket / disclosure-packet companions beside that rung \\
\bottomrule
\end{tabularx}
\caption{A forced public answer-escalation-ladder matrix. This is the smallest answer-side public answer later terse papers can quote directly when the live issue is whether one maintained reveal rung still forces the same widening target, rather than merely which ladder exists in the abstract.}
\label{tab:forced-public-answer-escalation-ladder-matrix}
\end{table}

This matrix is intentionally non-decorative: it pins the actual worked ladder-plus-rung pairs that still force the same maintained widening targets, together with the exact conditions under which that forcing remains honest or goes stale, without silently re-owning the imported publication-profile, citation-docket, carrier-slot, audit-trail, or disclosure-packet objects that neighboring notes already own.\cite{series:workedexample}

\paragraph{One tiny answer-escalation-ladder drill.}
For the worked auditor compact-diff answer, the pair \nolinkurl{auditor_compact_diff_summary_answer_escalation_ladder} plus \nolinkurl{exact_location} still forces \nolinkurl{auditor_compact_diff_summary_answer_carrier_slot}. If the live question changes to ``walk every semantic piece leftward,'' the honest stop remains the same ladder but the rung changes to \nolinkurl{full_audit}, so the forced widening target becomes \nolinkurl{auditor_compact_diff_summary_answer_audit_trail}. If the live question instead stays on \nolinkurl{exact_location} but the emitted answer card, sentence-lock lineage, or imported citation-docket / disclosure-packet companion agreement changes, the matrix has gone stale and a later terse paper must reopen the exact escalation-ladder owner rather than keep quoting the old carrier-slot rung. Only once the reveal rung itself is fixed and the real question becomes which already-maintained object should answer one concrete referee follow-up should it leave this note for Synthesis~54.\cite{series:workedexample}

\begin{thebibliography}{99}\small

\bibitem{series:challengeanswercitationdockets}
Anonymous.
\newblock Challenge Answer Citation Dockets for Successor Maintenance: Smallest-Sufficient Late-Stage Citation Choices from Capsules to Audit Trails.
\newblock Series synthesis note (Synthesis~50), 2026. (This archive.)

\bibitem{series:challengeanswerpublicationprofiles}
Anonymous.
\newblock Challenge Answer Publication Profiles for Successor Maintenance: Budgeted Reuse Modes from Answer Cards, Sentence Locks, and Citation Dockets.
\newblock Series synthesis note (Synthesis~51), 2026. (This archive.)

\bibitem{series:challengeanswerdisclosurepackets}
Anonymous.
\newblock Challenge Answer Disclosure Packets for Successor Maintenance: Public-on-Request Audit Slices from Publication Profiles, Audit Trails, and Support Manifests.
\newblock Series synthesis note (Synthesis~52), 2026. (This archive.)

\bibitem{series:challengeansweraudittrails}
Anonymous.
\newblock Challenge Answer Audit Trails for Successor Maintenance: Leftward Reader Walks from Exact Carrier Slots to Narrow Terminal Owners.
\newblock Series synthesis note (Synthesis~49), 2026. (This archive.)

\bibitem{series:challengeanswercarrierslots}
Anonymous.
\newblock Challenge Answer Carrier Slots for Successor Maintenance: Exact Shipped Fields and Reusable Clauses from Answer Cards.
\newblock Series synthesis note (Synthesis~48), 2026. (This archive.)

\bibitem{series:downstreamnormalforms}
Anonymous.
\newblock Downstream Normal Forms for Successor Maintenance: Typed Resolution and Exact Surface Assembly.
\newblock Series synthesis note (Synthesis~31), 2026. (This archive.)

\bibitem{series:releasespine}
Anonymous.
\newblock Release Spine Contracts for Successor Maintenance: Stable Stages and Stage-Local Handoffs.
\newblock Series synthesis note (Synthesis~32), 2026. (This archive.)

\bibitem{series:workedexample}
Anonymous.
\newblock Worked Example: Receipt Interlock for an Anonymous-DHT Reader-Privacy Claim.
\newblock Series synthesis note (Synthesis~17), 2026. (This archive.)

\bibitem{series:challengeanswerreviewmenus}
Anonymous.
\newblock Challenge Answer Review Menus for Successor Maintenance: Smallest-Sufficient Referee Handoffs from Dockets, Profiles, Packets, and Ladders.
\newblock Series synthesis note (Synthesis~54), 2026. (This archive.)


\bibitem{series:pairedterminalcitationdiscipline}
Anonymous.
\newblock Paired Terminal Citation Discipline for Review Spines: Stable Checked-Answer Stops and Adjacent Request-Side Capstones.
\newblock Series synthesis note (Synthesis~74), 2026. (This archive.)

\bibitem{series:pairedterminalcutoverrules}
Anonymous.
\newblock Paired Terminal Cutover Rules for Review Spines: When to Stop, Widen, or Fail Closed.
\newblock Series synthesis note (Synthesis~75), 2026. (This archive.)

\end{thebibliography}
\end{document}
